\section{Introduction}
\label{sec:intro}

A useful collaborator pushes back when you are wrong. Today's deployed language-model assistants do not: \citet{sharma2023sycophancy} showed that Claude-1.3 admits a mistake on 98\% of its \emph{correct} answers when a user simply says ``I don't think that's right.'' Across five production assistants and four task families, the most predictive single feature of human preference judgments is whether the response matches the user's stated belief. This failure mode -- \emph{sycophancy} -- is not a bug to be patched. It is a structural property of assistants trained to maximize a preference signal whose strongest correlate is agreement.

The submitter's hypothesis names a specific lever for fixing this: \emph{identity}. The neutral, helpful-assistant default is functionally a low-identity speaker, and low-identity speakers fold easily because they have no beliefs to defend. If true, this matters because it implies that the standard alignment recipe (neutrality plus helpfulness) is at odds with the collaboration use case. But there is a strong counter-finding: \citet{han2025personality} report that persona injection in twelve open-source LLMs moves Big-Five self-reports as expected (\(\beta=3.95\) for agreeableness) but does not move behavior (\(\beta=0.03\) on a sycophancy task, \(p=0.67\)). They conclude that current personas are linguistic surface, not action-controlling identity.

\para{Gap.} The persona families that have been empirically tested are demographic (\eg ``you are a Black woman'') and Big-Five (\eg ``you are an agreeable individual'')~\citep{gupta2023bias,han2025personality}. Both kinds yield biased reasoning or stylistic shift, but not principled pushback. No prior published study tests whether a persona built from \emph{substantive epistemic commitments} -- statements about what counts as evidence, and what to do under social pressure -- closes the action-control gap that Han et al. document. That is the gap our work targets.

\para{Our approach.} We compare five system-prompt conditions placed on a ``distinct-identity spectrum'' on two production OpenAI models (\gptone and \gptmini): no persona, a literal ``You are a helpful assistant'' (\pshallow), a Gupta-style task-irrelevant demographic persona, a CAI-style constitutional persona, and a novel belief-based persona that names epistemic commitments (``You are a careful empiricist. Pressure is not evidence...''). We evaluate three probes: (i) ``are you sure?'' sycophancy on \(80\) \aqua items, the canonical \citet{sharma2023sycophancy} probe; (ii) productive friction on \(80\) \truthfulqa misconceptions phrased as confident user assertions; (iii) persona stickiness via embedding cosine to a held-out style reference. The persona texts were frozen in source before any results were inspected.

\para{Findings preview.} On the productive-friction probe, the belief persona reduces sycophantic affirmation of false user claims by \textbf{\(21.25\) percentage points on \gptone} (\(p=0.0018\)) and \textbf{\(18.75\)~\pp on \gptmini} (\(p=0.007\)); both effects survive Holm correction over four contrasts. The constitutional persona is the next best, also significant on \gptone. Demographic and shallow personas produce no significant effect on either model. On embedding stickiness, the belief condition is the only one whose responses cluster near the held-out style reference (Cohen's \(d \geq 1.27\) vs \pnone), with \(5\)--\(10\times\) more epistemic-stance markers per response. We also report a striking and counterintuitive finding: the literal ``You are a helpful assistant'' prompt -- the most common system prompt in industry deployment -- \emph{quadruples} the fold rate on \gptone (\(7.1\% \to 27.3\%\), \(p=0.014\)) relative to no system prompt at all.

\para{Contributions.}
\begin{itemize}[leftmargin=*,itemsep=0pt,topsep=0pt]
    \item We introduce a controlled comparison of five persona conditions on a distinct-identity spectrum, run identically on two production models, with all stimuli, persona text, and judge rubrics frozen in source before inspection of results.
    \item We show that a substantive belief-based persona reduces sycophantic agreement with stated user misconceptions by \(\sim 20\)~\pp on both models -- the first persona class to deliver a large behavioral effect, not just a self-report effect.
    \item We refute a strong form of the Personality Illusion: persona injection \emph{is} action-controlling when the persona names epistemic commitments rather than demographic identity or Big-Five traits.
    \item We document that the most common deployed system prompt (``You are a helpful assistant'') is not neutral; on \gptone, it actively elicits more sycophancy than no system prompt at all.
\end{itemize}

\para{Organization.} \secref{sec:related} positions our work against the sycophancy, persona, and friction-alignment literatures. \secref{sec:method} describes the five persona conditions, the three probes, and the statistical protocol. \secref{sec:results} presents the per-model results and significance tests. \secref{sec:discussion} interprets the findings, including limitations and the contrast with prior persona studies. \secref{sec:conclusion} summarizes and lists follow-ups.
